\section{Discussion}
\label{sec:discussion}

The computer systems literature is rich with coordination mechanisms for managing concurrency
both within and across machines. When the CALM Conjecture was first being considered,
it addressed a different question: when are these mechanisms required, and why? 
The original CALM Conjecture built on emerging work in declarative
networking \jmh{cite Loo} for intuition. Inspired by Codd's tradition of
separating specification from implementation \jmh{cite Codd}, declarative
networking showed that the desired effects of protocols could be written as
queries, and subsequently analyzed and optimized into a variety of mechanisms.

There was a related friction at the time in the literature on parallel query processing. The
widely-cited Google MapReduce paper \jmh{cite} departed from prior work on parallel
query processing~\jmh{cite WilshutApers, Exchange}, conflating data redistribution
(shuffle) with barrier coordination required for a non-monotone reduce operator.
This poor separation of concerns required monotone operators like join to block
in the MapReduce model unnecessarily.  This conflation leaked into theory as
well: metrics of ``rounds'' of BSP or MapReduce effectively charge a complexity
cost to incidental uses of barriers that similarly are not intrinsic to program
semantics \jmh{cite Afrati/Ullman, BSP, who else}.

The CALM Conjecture hypothesized that the database community's tradition of more
judicious use of barriers could be generalized: coordination---of any kind---is
required if and only if the program is non-monotone. The CALM Theorem of Ameloot
et al.\ proved it. Unfortunately the impact of the conjecture and its proof was
limited by its embrace of declarative languages, which to date have still not
become mainstream (more on this below.)

This paper can be seen as moving one level higher in its declarations, from
programs made of queries to a single virtual table---a view. For our analysis,
what is important about a specification is only the encapsulated view mapping
its base input to its final output. By eliding the query logic entirely, we not
only reduce complexity, we free ourselves from the syntax and semantics of any
particular language, providing a declarative view compatible with any
programming framework.

By raising the abstraction, the focus of what to specify naturally broadens as
well. Prior work focused on implementing complex programs declaratively
\jmh{cite loo-sigmod, chu, P2, Bloom}. This was interesting ``declarative
engineering'' with concrete benefits in analysis and optimization.  The model
here naturally expands the scope, naturally modeling abstract families of
invariants like CAP, isolation levels or consistency requirements that have no
specific application or protocol in mind.

Lifting the abstraction improves the power of analysis as well. Prior attempts
to do program analysis, e.g.\ in languages like Bloom~\cite{alvaro2011bloom},
became enmeshed in the details of the ``declarative'' syntax of full
implementations. Bloom could conservatively identify locations in programs where
coordination protocols would likely be required based on syntax: the use of
potentially non-monotone constructs. But adding coordination code in Bloom does
not remove those constructs---it actually adds more. Our composition results
here (\S\ref{sec:composition}) benefit from the absence of programs. The
resaulting analysis is exact rather than conservative \jmh{soften that claim!},
and via proper coordination shows how the non-monotonicity of a coordination
protocol frees its consumers from the need to coordinate further.

Underlying all this discussion is a refined question of what a ``declarative
specification'' should capture---a reassessment of Codd's agenda. This is
timely. Automated code generation in imperative languages is now a reality, but
the ensuing large volume of code is often rife with bugs. This has led to more
widespread interest in formal specification languages like Lean \jmh{cite}, and
a renewed opportunity for logic, dataflow or functional programming languages to
capture the attention of developers and their coding agents. Still, these
languages are typically used by developers to specify protocols and programs at
a significant level of detail, as in the declarative networking work. The
question of ``how declarative should my specification be'' is increasingly
topical. This paper points to an approach that specifies at a higher level,
which may offer new opportunities for both analyses like ours, as well as formal
program synthesis and optimization.
